% BHU PYQ -- Manually transcribed & error-corrected
% Paper: BCF-213 : Corporate Taxation
% Exam: B.Com. (Hons.) FMM Semester III Examination, 2022-23

\documentclass[11pt,a4paper]{article}
\usepackage[a4paper,top=2.5cm,bottom=2.5cm,left=2.8cm,right=2.8cm,headheight=14pt]{geometry}
\usepackage{amsmath,amssymb}
\usepackage[T1]{fontenc}
\usepackage[utf8]{inputenc}
\usepackage{lmodern,microtype,fancyhdr,enumitem,hyperref,lastpage,parskip,booktabs}

\hypersetup{
    colorlinks=true,
    urlcolor=black,
    linkcolor=black,
    pdftitle={BCF-213: Corporate Taxation -- BHU B.Com. (Hons.) FMM Sem III 2022-23}
}

\pagestyle{fancy}
\fancyhf{}
\renewcommand{\headrulewidth}{0.4pt}
\renewcommand{\footrulewidth}{0.4pt}
\fancyhead[L]{\small   \textit{Banaras Hindu University}}
\fancyhead[R]{\small   \textit{BCF-213: Corporate Taxation}}
\fancyfoot[C]{\small Page  \thepage\ of \pageref{LastPage}}


\newlist{parts}{enumerate}{2}
\setlist[parts,1]{label=(\alph*),leftmargin=*,itemsep=5pt,topsep=3pt}
\setlist[parts,2]{label=(\roman*),leftmargin=*,itemsep=3pt,topsep=2pt}

\newcommand{\pts}[1]{\hfill   \textit{[#1]}}


\begin{document}


\begin{center}
  {\large\bfseries BANARAS HINDU UNIVERSITY}\\[2pt]
  {\normalsize B.Com. (Hons.) FMM --- Semester III Examination, 2022-23}\\[6pt]
  \rule{0.8\linewidth}{0.5pt}\\[6pt]
  {\large\bfseries Paper No. BCF-213: Corporate Taxation}\\[4pt]
  {\normalsize Time: 3 Hours\hfill Maximum Marks: 70}
\end{center}

\rule{\linewidth}{0.4pt}

\smallskip



\noindent   \textit{Instructions: Answer all questions. All questions carry equal marks (14 marks each).}

\bigskip




\noindent   \textbf{Question 1.} \\
Mr. Ram furnishes the following particulars of his income earned during the previous year relevant to the assessment year 2022-23. Compute his Gross Total Income if he is: (a) Ordinary Resident, (b) Not-Ordinary Resident, (c) Non-Resident.


\begin{table}[h!]
  \centering
  \small
  
\begin{tabular}{clr}
     oprule
   \textbf{Sl. No.} &   \textbf{Particulars} &   \textbf{Amount (Rs.)} \\
    \midrule
    (i)   & Interest on Savings Bank Deposit in Corporation Bank, Delhi & 12,000 \\
    (ii)  & Income from agriculture in Africa, invested in Russia & 5,000 \\
    (iii) & Dividends received in USA from English Company, out of which Rs. 2,000 were remitted to India & 12,000 \\
    (iv)  & Salary drawn for two months for working in Indian Embassy's Office in Australia & \\
          & and salary received there & 48,000 \\
    (v)   & Income from house property (the building is situated in Iraq), out of which Rs. 20,000 & \\
          & was deposited in a bank in Iraq and the balance remitted to India & 25,000 \\
    (vi)  & Pension received in Belgium for services rendered in India with a limited liability company & 10,000 \\
    (vii) & Interest paid by an Indian company but received in London & 2,00,000 \\
    ottomrule
  \end{tabular}
\end{table}\pts{14}

\centerline{   \textbf{OR}}

What are the different categories into which the assessees are divided with regard to residence? Give a brief account of each of them.\pts{14}

\medskip\rule{\linewidth}{0.2pt}

\pagebreak




\noindent   \textbf{Question 2.} \\
The following is the Receipts \& Payments Account of Mr. `GC', a practicing Chartered Accountant for the year ended 31-03-2022:


\begin{table}[h!]
  \centering
  \small
  
\begin{tabular}{lrlr}
     oprule
   \textbf{Receipts} &   \textbf{Amount (Rs.)} &   \textbf{Payments} &   \textbf{Amount (Rs.)} \\
    \midrule
    Audit fees & 19,210 & Office expenses & 10,000 \\
    Consultation fee & 10,000 & Office rent & 5,000 \\
    Appellate Tribunal appearance fee & 15,000 & Salaries and wages & 12,050 \\
    Miscellaneous fee & 20,000 & Printing and stationery & 1,000 \\
    Interest on Government Securities & 10,000 & Subscription to CA Institute & 3,000 \\
    Rent received & 10,000 & Purchase of books for professional purposes & 1,300 \\
    Presents from client & 10,000 & Travelling expenses & 5,800 \\
          &        & Interest on bank loan & 3,000 \\
          &        & Donation to National Defense Fund & 5,000 \\
    ottomrule
  \end{tabular}
\end{table}




\noindent   \textbf{Additional Information:}

\begin{enumerate}[label=(\alph*)]
  \item Loan from bank was taken for the construction of the house in which he lives. The municipal value of this house is Rs. 8,000 and local taxes Rs. 800 p.a.
  \item 1/4th of travelling expenses are not allowable.
\end{enumerate}
Compute Professional Income of Mr. GC for the assessment year 2022-23.\pts{14}

\centerline{   \textbf{OR}}

What is Book Profit? Under Section 40(b), what items are disallowed as deduction while computing firm's income from business or profession? Explain fully.\pts{14}

\medskip\rule{\linewidth}{0.2pt}




\noindent   \textbf{Question 3.} \\
X, Y and Z, the partners share profits and losses in the ratio of 1:2:3. Net profit as per Profit \& Loss Account is Rs. 1,31,800 for the previous year 2021-22. In computing the total income of Rs. 1,31,800, the following have been debited to the P \& L Account:

\begin{enumerate}
  \item Salaries of Rs. 1,30,000.
  \item Interest on Capital calculated @ 20\% of Rs. 3,500, Rs. 14,000 and Rs. 10,500 to X, Y and Z respectively.
  \item Bonus to Z Rs. 15,000.
  \item Commission of Rs. 5,000, Rs. 12,500 and Rs. 17,500 to X, Y and Z respectively.
\end{enumerate}
Z had borrowed capital for his investment in the firm and had paid interest of Rs. 7,500 separately to the lender.

Compute the Total Income of the firm and Taxable Income of the three partners in the firm. All are working partners. The firm fulfils the conditions of Section 184 of the Income-tax Act, 1961.\pts{14}

\centerline{   \textbf{OR}}

Explain clearly the deductions expressly admissible and inadmissible in computing the Income from `Business and Profession' as per Income-tax Act, 1961.\pts{14}

\medskip\rule{\linewidth}{0.2pt}

\pagebreak




\noindent   \textbf{Question 4.} \\
The following are the incomes of Mr. R for the financial year 2022-23:
% NOTE: Professional income of 7,50,000 is printed in the paper, but to get the printed Gross Total Income of 11,85,000, 
% the professional income must mathematically be 14,50,000. We keep 7,50,000 as printed, with a LaTeX comment.

\begin{itemize}
  \item Taxable Income from Profession: Rs. 7,50,000 (Printed as such, mathematically requires 14,50,000)
  \item Taxable Income from House Property: Rs. $-$3,50,000
  \item Dividend from Indian Company: Rs. 50,000
  \item Interest on Savings Bank deposit: Rs. 25,000
  \item Interest on Post Office deposits: Rs. 10,000
  \item Gross Total Income: Rs. 11,85,000
\end{itemize}

Additional transactions:

\begin{enumerate}[label=(\alph*)]
  \item He has paid Life Insurance Premium: Rs. 70,000
  \item He has paid Health Insurance Premium by cheques: Rs. 10,000
  \item He has donated to PM Cares Fund by cheques: Rs. 15,000
\end{enumerate}
Calculate Advance Tax Installments u/s 211 of Income-tax Act, 1961.\pts{14}

\centerline{   \textbf{OR}}

Write short notes on the following:\pts{7+7}

\begin{parts}
  \item Unabsorbed depreciation
  \item Procedure of an appeal to the Commissioner (Appeals)
\end{parts}

\medskip\rule{\linewidth}{0.2pt}




\noindent   \textbf{Question 5.} \\
Mr. B is an employee of a company at Varanasi. He received the following during the financial year 2022-23:

\begin{itemize}
  \item Salary: Rs. 9,60,000
  \item Dearness allowance: Rs. 2,80,000
  \item Bonus: Rs. 80,000
  \item Children Education Allowance: Rs. 150 per month for one child
  \item Chargeable/Computed value of rent-free house: Rs. 1,56,000
  \item His annual contribution to Provident Fund: Rs. 96,000
  \item He has paid Life Insurance Premium: Rs. 12,000
  \item Donation to PM Cares Fund: Rs. 10,000
\end{itemize}
Compute his Income from Salary liable to deduction of tax at source and the amount of tax to be deducted.\pts{14}

\centerline{   \textbf{OR}}

What do you understand by `Advance Payment of Tax'? Explain the provisions of the Income-tax Act in this respect.\pts{14}

\vfill
\rule{\linewidth}{0.4pt}

\begin{center}\small
  \href{https://bhu-exam.vercel.app/}{Click here to visit our website}\\[4pt]
  \href{https://me-aryan.vercel.app/}{Click here to visit our contributor}\\[4pt]
  \href{https://quashan-me.vercel.app/}{Click here to visit our contributor}
\end{center}

\end{document}
